% ═══════════════════════════════════════════════════════════════
% LERNBLATT: Reaktionsgleichungen aufstellen und ausgleichen
% Thema 2.5 | Chemie Klasse 8 | Kompakt 2 Seiten
% ═══════════════════════════════════════════════════════════════

\documentclass[9pt, a4paper]{article}

\usepackage[utf8]{inputenc}
\usepackage[T1]{fontenc}
\usepackage[ngerman]{babel}

\usepackage{helvet}
\renewcommand{\familydefault}{\sfdefault}

\usepackage{microtype}
\usepackage[margin=1.2cm, top=1.4cm, bottom=1cm]{geometry}
\usepackage{enumitem}
\usepackage{tabularx}
\usepackage{booktabs}
\usepackage{array}
\usepackage{multicol}

\usepackage{amsmath, amssymb}
\usepackage{siunitx}
\sisetup{locale = DE, per-mode = symbol}

\usepackage{fancyhdr}
\usepackage{lastpage}

\usepackage{xcolor}
\usepackage{tcolorbox}
\tcbuselibrary{skins}

\usepackage{tikz}
\usetikzlibrary{shapes.geometric, arrows.meta, positioning}

\usepackage{qrcode}

% ═══════════════════════════════════════════════════════════════
% FARBEN
% ═══════════════════════════════════════════════════════════════

\definecolor{chemie}{HTML}{2a7a4b}
\definecolor{hellgrau}{HTML}{F5F5F5}
\definecolor{mittelgrau}{HTML}{888888}
\definecolor{dunkelgrau}{HTML}{444444}
\definecolor{metall}{HTML}{2c5aa0}
\definecolor{sauerstoff}{HTML}{c62828}

\colorlet{fachfarbe}{chemie}
\colorlet{fachfarbe-hell}{fachfarbe!15}

% ═══════════════════════════════════════════════════════════════
% KOPF- UND FUSSZEILE
% ═══════════════════════════════════════════════════════════════

\pagestyle{fancy}
\fancyhf{}
\renewcommand{\headrulewidth}{0pt}
\renewcommand{\footrulewidth}{0pt}
\fancyhead[L]{\small\textcolor{fachfarbe}{\textbf{Chemie}} | Thema 2.5}
\fancyhead[R]{\small Klasse 8}
\fancyfoot[C]{\small\textcolor{mittelgrau}{\thepage}}

% ═══════════════════════════════════════════════════════════════
% BOXEN
% ═══════════════════════════════════════════════════════════════

\newtcolorbox{formelkasten}{
    colback=fachfarbe-hell,
    colframe=fachfarbe,
    boxrule=1pt,
    arc=0pt,
    left=8pt, right=8pt, top=6pt, bottom=6pt,
    before skip=4pt, after skip=4pt
}

\newtcolorbox{merkkasten}{
    colback=hellgrau,
    colframe=hellgrau,
    boxrule=0pt,
    arc=0pt,
    left=6pt, right=6pt, top=5pt, bottom=4pt,
    borderline west={2pt}{0pt}{fachfarbe},
    before skip=4pt, after skip=4pt
}

\newtcolorbox{warnkasten}{
    colback=sauerstoff!10,
    colframe=sauerstoff,
    boxrule=1pt,
    arc=0pt,
    left=6pt, right=6pt, top=5pt, bottom=4pt,
    before skip=4pt, after skip=4pt
}

\newtcolorbox{hinweiskasten}{
    colback=metall!10,
    colframe=metall,
    boxrule=1pt,
    arc=0pt,
    left=6pt, right=6pt, top=5pt, bottom=4pt,
    before skip=4pt, after skip=4pt
}

\newtcolorbox{algorithmusbox}[1][]{
    enhanced,
    colback=fachfarbe!8,
    colframe=fachfarbe,
    boxrule=1.5pt,
    arc=0pt,
    left=6pt, right=6pt, top=10pt, bottom=5pt,
    before skip=8pt, after skip=4pt,
    title={#1},
    fonttitle=\bfseries\small,
    coltitle=white,
    attach boxed title to top left={yshift=-3mm, xshift=4mm},
    boxed title style={colback=fachfarbe, arc=0pt, boxrule=0pt, left=4pt, right=4pt}
}

% ═══════════════════════════════════════════════════════════════
% BEFEHLE
% ═══════════════════════════════════════════════════════════════

\newcommand{\seitentitel}[1]{%
    \noindent\textcolor{fachfarbe}{\rule{\textwidth}{1.5pt}}\\[-0.1em]
    {\Large\bfseries\textcolor{fachfarbe}{#1}}\\[-0.3em]
    \textcolor{fachfarbe}{\rule{\textwidth}{0.4pt}}
    \vspace{0.3em}
}

\newcommand{\abschnitt}[1]{%
    \vspace{6pt}%
    \noindent{\normalsize\bfseries\textcolor{fachfarbe}{#1}}%
    \vspace{2pt}%
    \nopagebreak[4]%
}

\setlength{\parindent}{0pt}
\setlength{\parskip}{2pt}
\setlength{\columnsep}{0.8cm}

% ═══════════════════════════════════════════════════════════════
% DOKUMENT
% ═══════════════════════════════════════════════════════════════

\begin{document}

% ═══════════════════════════════════════════════════════════════
% SEITE 1
% ═══════════════════════════════════════════════════════════════

\seitentitel{Reaktionsgleichungen aufstellen und ausgleichen}

\begin{multicols}{2}

\abschnitt{Warum Reaktionsgleichungen?}

Reaktionsgleichungen sind die \textbf{Rezepte der Chemie}. Sie zeigen die genauen Mengenverhältnisse.

\textbf{Beispiel Airbag:} Die exakte Menge Natriumazid produziert genau 67\,L Stickstoff in 30\,ms – nicht mehr, nicht weniger!

\abschnitt{Wertigkeit – Die „Händezahl"}

Die \textbf{Wertigkeit} gibt an, wie viele Bindungen ein Atom eingehen kann.

\vspace{4pt}
\begin{center}
\small
\begin{tabular}{|c|c|c|c|c|c|c|c|}
\hline
\textbf{Na} & \textbf{Mg} & \textbf{Al} & \textbf{Fe} & \textbf{Cu} & \textbf{Zn} & \textbf{Ca} & \textbf{O} \\
\hline
I & II & III & II/III & I/II & II & II & II \\
\hline
\end{tabular}
\end{center}

\vspace{4pt}
\textbf{Vorstellung:} Mg hat 2 „Hände", Al hat 3 „Hände", O hat 2 „Hände".

\columnbreak

\abschnitt{Die Kreuzregel}

Bei \textbf{unterschiedlichen} Wertigkeiten:

\begin{formelkasten}
\centering
Wertigkeit von A → Index von B\\
Wertigkeit von B → Index von A
\end{formelkasten}

\textbf{Beispiel:}\\
\textcolor{metall}{Al}(III) + \textcolor{sauerstoff}{O}(II) → \textcolor{metall}{Al}\textsubscript{\textcolor{sauerstoff}{2}}\textcolor{sauerstoff}{O}\textsubscript{\textcolor{metall}{3}}

\vspace{4pt}
\small
\begin{tabular}{@{}ll@{}}
Mg(II) + O(II) & → MgO (1:1) \\
Al(III) + O(II) & → Al\textsubscript{2}O\textsubscript{3} \\
Fe(III) + O(II) & → Fe\textsubscript{2}O\textsubscript{3} \\
Na(I) + O(II) & → Na\textsubscript{2}O \\
\end{tabular}

\end{multicols}

\vspace{-4pt}

\begin{hinweiskasten}
\small
\textbf{Warum bilden sich diese Formeln?} Bei Metalloxiden gibt das Metall Elektronen \textbf{ab}, Sauerstoff \textbf{nimmt} sie auf. So entstehen geladene Teilchen (\textbf{Ionen}), die sich anziehen. Die Wertigkeit zeigt, wie viele Elektronen abgegeben bzw. aufgenommen werden.
\end{hinweiskasten}

\vspace{-2pt}

% ═══════════════════════════════════════════════════════════════
% DIE ZWEI ALGORITHMEN
% ═══════════════════════════════════════════════════════════════

\begin{algorithmusbox}[Algorithmus 1: Verhältnisformel aufstellen]
\small
\begin{tabular}{@{}cl}
\textbf{1.} & Elementsymbole aufschreiben: \hfill Li \quad O \\
\textbf{2.} & Wertigkeiten ablesen: \hfill I \quad II \\
\textbf{3.} & Kreuzregel: Wertigkeit A → Index B \hfill → \textbf{Li\textsubscript{2}O} \\
\end{tabular}
\end{algorithmusbox}

\vspace{-4pt}

\begin{algorithmusbox}[Algorithmus 2: Reaktionsgleichung aufstellen und ausgleichen]
\small
\begin{tabular}{@{}cl}
\textbf{1.} & Wortgleichung: Lithium + Sauerstoff → Lithiumoxid \\
\textbf{2.} & Formelgleichung: Li + O\textsubscript{2} → Li\textsubscript{2}O \quad {\scriptsize(Achtung: Sauerstoff ist O\textsubscript{2}!)} \\
\textbf{3.} & Atome zählen: Links: Li=1, O=2 \quad Rechts: Li=2, O=1 \\
\textbf{4.} & Vorfaktoren ausgleichen: \\
& \quad a) Produkt verdoppeln, bis O-Zahl gerade: Li\textsubscript{2}O hat 1 O $\rightarrow$ 2 Li\textsubscript{2}O hat 2 O \\
& \quad b) O\textsubscript{2}-Anzahl = O rechts $\div$ 2 = 1 \\
& \quad c) Metall-Anzahl = Metall-Atome rechts = 4 \\
& \quad $\Rightarrow$ \textbf{4} Li + \textbf{1} O\textsubscript{2} $\rightarrow$ \textbf{2} Li\textsubscript{2}O \\
& \quad d) \textbf{Kontrolle:} Li: 4=4 \checkmark \quad O: 2=2 \checkmark \\
\end{tabular}
\end{algorithmusbox}

% ═══════════════════════════════════════════════════════════════
% SEITE 2
% ═══════════════════════════════════════════════════════════════

\newpage

\seitentitel{Reaktionsgleichungen aufstellen und ausgleichen}

\begin{multicols}{2}

\abschnitt{Häufige Fehler vermeiden}

\begin{warnkasten}
\textbf{Fehler 1: O statt O\textsubscript{2}}\\
Sauerstoff kommt in der Natur als \textbf{O\textsubscript{2}} vor!\\
\textcolor{sauerstoff}{$\times$} Mg + O $\rightarrow$ MgO\\
\textcolor{fachfarbe}{$\checkmark$} Mg + O\textsubscript{2} $\rightarrow$ MgO
\end{warnkasten}

\vspace{4pt}

\begin{warnkasten}
\textbf{Fehler 2: Index statt Vorfaktor}\\
\textbf{Index} (tiefgestellt) ändert den \textbf{Stoff}!\\
MgO $\neq$ MgO\textsubscript{2} (anderer Stoff!)\\[0.3em]
\textbf{Vorfaktor} (vor der Formel) ändert nur die \textbf{Menge}:\\
2 MgO = zwei Mal MgO
\end{warnkasten}

\columnbreak

\abschnitt{Ausgeglichene Gleichungen}

\small
\begin{tabular}{@{}p{5cm}l@{}}
\textbf{Magnesium:} & \textbf{2} Mg + O\textsubscript{2} → \textbf{2} MgO \\
& {\scriptsize(2 Mg, 2 O = 2 Mg, 2 O)} \\[0.5em]
\textbf{Aluminium:} & \textbf{4} Al + \textbf{3} O\textsubscript{2} → \textbf{2} Al\textsubscript{2}O\textsubscript{3} \\
& {\scriptsize(4 Al, 6 O = 4 Al, 6 O)} \\[0.5em]
\textbf{Eisen:} & \textbf{4} Fe + \textbf{3} O\textsubscript{2} → \textbf{2} Fe\textsubscript{2}O\textsubscript{3} \\
& {\scriptsize(4 Fe, 6 O = 4 Fe, 6 O)} \\[0.5em]
\textbf{Kupfer:} & \textbf{2} Cu + O\textsubscript{2} → \textbf{2} CuO \\
& {\scriptsize(2 Cu, 2 O = 2 Cu, 2 O)} \\[0.5em]
\textbf{Natrium:} & \textbf{4} Na + O\textsubscript{2} → \textbf{2} Na\textsubscript{2}O \\
& {\scriptsize(4 Na, 2 O = 4 Na, 2 O)} \\
\end{tabular}

\end{multicols}

\vspace{4pt}

\abschnitt{Tipp: kgV finden}

Bei unterschiedlichen O-Zahlen links und rechts: Finde das \textbf{kleinste gemeinsame Vielfache} (kgV).

\vspace{4pt}
\small
\begin{tabular}{|l|c|c|c|l|}
\hline
\textbf{Gleichung} & \textbf{O links} & \textbf{O rechts} & \textbf{kgV} & \textbf{Lösung} \\
\hline
Mg + O\textsubscript{2} → MgO & 2 & 1 & 2 & 2 MgO, also 2 Mg \\
\hline
Al + O\textsubscript{2} → Al\textsubscript{2}O\textsubscript{3} & 2 & 3 & 6 & 3 O\textsubscript{2}, 2 Al\textsubscript{2}O\textsubscript{3}, 4 Al \\
\hline
Na + O\textsubscript{2} → Na\textsubscript{2}O & 2 & 1 & 2 & 2 Na\textsubscript{2}O, also 4 Na \\
\hline
\end{tabular}

\vspace{8pt}

\begin{merkkasten}
\textbf{Das Wichtigste:}
\begin{itemize}[leftmargin=1.2em, itemsep=0pt, topsep=2pt]
    \item \textbf{Wertigkeit} = Anzahl der möglichen Bindungen („Händezahl") = Elektronen, die abgegeben/aufgenommen werden
    \item \textbf{Kreuzregel:} Wertigkeit A → Index B, Wertigkeit B → Index A
    \item Sauerstoff ist immer \textbf{O\textsubscript{2}} (zweiatomig)
    \item \textbf{Index} ändert den Stoff, \textbf{Vorfaktor} ändert die Menge
    \item \textbf{Massenerhaltung:} Atome links = Atome rechts
\end{itemize}
\end{merkkasten}

% ═══════════════════════════════════════════════════════════════
% QR-CODES
% ═══════════════════════════════════════════════════════════════

\vfill

\begin{center}
\begin{tabular}{cc}
\begin{minipage}{2.5cm}
\centering
\qrcode[height=1.5cm]{https://devbox42.github.io/sciencesim/chemie/kl08/02-reaktionsgleichungen/LP-02-reaktionsgleichungen.html}\\
{\tiny Lernpfad}
\end{minipage}
&
\begin{minipage}{2.5cm}
\centering
\qrcode[height=1.5cm]{https://devbox42.github.io/sciencesim/chemie/kl08/02-reaktionsgleichungen/MT-02-reaktionsgleichungen.html}\\
{\tiny Minitest}
\end{minipage}
\end{tabular}
\end{center}

\end{document}
